\chapter{Modèles séquentiels sur les carnets d'ordres}\label{ch:ml:sequence-models-on-order-books}

Un article très cité rapporte qu'un réseau convolutif-récurrent lisant les états bruts d'un carnet d'ordres à dix niveaux
surpasse tous les algorithmes existants sur un jeu de données de référence de carnets d'ordres, et offre une précision stable hors échantillon (out-of-sample) sur une
année de cotations de la London Stock Exchange. Sur un ordinateur portable,
sur quelques heures d'un carnet simulé, avec des étiquettes construites honnêtement et chaque session de test tenue à l'écart de l'entraînement, le meilleur
modèle séquentiel atteint un coefficient d'information (IC) de 0.43, et une \omterm{def:ml:linear-and-regularised-baselines:logistic}{régression logistique} sur dix variables construites à la main (le
déséquilibre du flux d'ordres, le déséquilibre des files, le flux récent et les mouvements) atteint 0.54. L'article n'a pas tort sur ses données~; le résultat est que
les carnets bruts sont une entrée difficile, les modèles séquentiels en exigent beaucoup, et l'affirmation «~l'apprentissage profond trouve
lui-même les variables~» a un coût en données qu'il faut payer en sessions. Ce chapitre construit les quatre architectures séquentielles de référence,
les entraîne sur des fenêtres de carnet d'ordres et mesure ce qu'elles apportent.

\section{Flux de carnets d'ordres en entrée}

Les données du chapitre sont des sessions de dix minutes de \texttt{firm.tape} (livre~7, chapitre~2), chacune avec sa propre graine~: huit pour
l'entraînement, trois pour l'\omterm{def:ml:trees-and-boosting:early}{arrêt précoce} (early stopping), huit pour le test, et huit autres pour doubler plus tard l'\omterm{def:ml:why-financial-machine-learning-is-different:sets}{ensemble d'entraînement}. Toutes les
demi-secondes, le carnet est reconstruit à partir des messages et résumé niveau par niveau~: pour chacun des dix meilleurs niveaux de chaque côté,
la distance du prix au mid, en ticks, et le logarithme de la taille en lots, soit quarante nombres par instantané. Une
entrée est une fenêtre des 20 derniers instantanés (dix secondes)~; son étiquette est la direction de la moyenne du mid sur les dix
secondes suivantes, par rapport au mid courant~: baisse, stable ou hausse, avec une zone morte d'un quart de tick (la \cref{fig:ml:lobseq:window}).
Les fenêtres ne franchissent jamais les sessions, et la mise à l'échelle des quarante colonnes n'est calibrée que sur les sessions d'entraînement.

\begin{omfigure}
\begin{tikzpicture}[x=0.40cm,y=0.40cm,font=\footnotesize]
\foreach \t in {0,...,19} { \foreach \l in {0,...,9} {
  \fill[omDef!15] (\t,\l) rectangle ++(0.85,0.85); } }
\draw[omDef,thick] (-0.2,-0.2) rectangle (19.95,10.05);
\node[anchor=east,align=right] at (-0.5,5) {10 niveaux $\times$ 4\\(distance, log de la taille,\\bid et ask)};
\node at (10,-1.1) {20 instantanés, un par demi-seconde (dix secondes)};
\draw[->,thick] (20.4,5) -- (23.3,5) node[midway,above] {modèle};
\node[draw=omThm,fill=omThm!10,rounded corners=2pt,align=center,anchor=west] at (23.5,5) {$\P(\text{baisse})$, $\P(\text{stable})$,\\ $\P(\text{hausse})$};
\draw[omProp,thick,->] (19.9,11.2) -- (26.5,11.2);
\node[omProp,anchor=south] at (23.2,11.3) {10 s suivantes~: mid moyen $-$ mid};
\draw[omProp,dashed] (19.9,10.2) -- (19.9,11.6);
\node[anchor=south] at (10,10.2) {instant de décision $t$ à la fin de la fenêtre};
\end{tikzpicture}

\omcaption{Une fenêtre d'entrée et son étiquette. L'étiquette compare la moyenne du mid sur les dix secondes suivantes avec le mid à
l'instant de décision et reste stable à un quart de tick près~; le lissage suit la littérature DeepLOB et rend l'étiquette
plus facile à prédire que le mouvement suivant lui-même.}\label{fig:ml:lobseq:window}
\end{omfigure}

Les classes sont déséquilibrées~: 73~\% des fenêtres de test sont stables, 15~\% en baisse et 13~\% en hausse. Prédire toujours «~stable~» est
juste dans 73~\% des cas, si bien que la précision est un mauvais score~; le chapitre la rapporte pour la comparer à la littérature et classe
les modèles selon le coefficient d'information de $\P(\text{hausse}) - \P(\text{baisse})$ avec la variation future, session par
session.

\section{Modèles convolutifs et récurrents}

\begin{definition}[Réseau de neurones convolutif]\label{def:ml:sequence-models-on-order-books:cnn}
Un \emph{réseau de neurones convolutif}\index{réseau de neurones convolutif} applique le même petit filtre à chaque
position de son entrée~: une couche convolutive unidimensionnelle envoie une séquence $z_1,\dots,z_T$ de vecteurs à $d$ composantes vers
$u_t = \mathrm{act}(\sum_{j=0}^{k-1}B_jz_{t+j-\lfloor k/2\rfloor} + c)$, avec $k$ matrices de filtre $B_j$ communes à tous les $t$.
\end{definition}

\begin{definition}[Réseau de neurones récurrent, mémoire à long et court terme]\label{def:ml:sequence-models-on-order-books:rnn}
Un \emph{réseau de neurones récurrent}\index{réseau de neurones récurrent} lit une séquence pas à pas et porte un état caché,
$h_t = \mathrm{act}(A x_t + C h_{t-1} + c)$. Un réseau à \emph{mémoire à long et court terme}\index{mémoire à long et court terme}
(LSTM) ajoute un état de cellule mis à jour par des portes (entrée, oubli et sortie, chacune une sigmoïde de $x_t$ et $h_{t-1}$), ce qui
laisse les gradients se propager sur de longues séquences (Hochreiter et Schmidhuber, 1997).
\end{definition}

Le CNN du chapitre comporte deux couches convolutives de 32 filtres de largeur cinq, en fait la moyenne sur la fenêtre et classe~;
le LSTM compte 32 unités et classe à partir de son dernier état. Les deux lisent les quarante colonnes de chaque instantané de la fenêtre.

\section{Convolutions temporelles et attention}

\begin{definition}[Convolution causale, réseau convolutif temporel]\label{def:ml:sequence-models-on-order-books:tcn}
Une \emph{convolution causale}\index{convolution causale} n'utilise que les positions courante et antérieures~:
$u_t = \sum_{j=0}^{k-1}B_jz_{t - jr}$, avec une dilatation $r$. Un
\emph{réseau convolutif temporel}\index{réseau convolutif temporel} (TCN) empile des convolutions causales avec des dilatations $1, 2, 4,\dots$, de sorte que le champ réceptif croisse
exponentiellement avec la profondeur, généralement avec des connexions résiduelles (Bai, Kolter et Koltun, 2018).
\end{definition}

\begin{definition}[Mécanisme d'attention, transformer]\label{def:ml:sequence-models-on-order-books:attention}
Un \emph{mécanisme d'attention}\index{mécanisme d'attention} permet à chaque position d'une séquence de former sa sortie comme
moyenne pondérée de toutes les positions~: avec des requêtes $\mathsf Q = ZA_{\mathsf Q}$, des clés $\mathsf K = ZA_{\mathsf K}$ et des valeurs
$\mathsf V = ZA_{\mathsf V}$ calculées à partir de la séquence $Z$ ($T\times h$), la sortie est
$\mathrm{softmax}(\mathsf Q\mathsf K^\top/\sqrt h)\mathsf V$. Une couche \emph{transformer}\index{transformer} combine
plusieurs têtes d'attention, un réseau appliqué position par position, des connexions résiduelles et une \omterm{def:ml:neural-networks-for-noisy-tabular-data:dropout}{normalisation par couche} (layer normalisation)~; les positions sont codées
en ajoutant aux entrées des vecteurs appris ou fixes (Vaswani et al., 2017).
\end{definition}

Le TCN compte quatre blocs causaux (dilatations de 1 à 8, filtres de largeur trois, champ réceptif de 31 instantanés, plus que la
fenêtre) et classe à partir de la dernière position~; le modèle d'attention est une couche \omterm{def:ml:sequence-models-on-order-books:attention}{transformer} à deux têtes de largeur 32
et positions apprises, lu lui aussi à la dernière position.

\begin{omfigure}
\begin{tikzpicture}
\begin{axis}[width=10.5cm,height=5.4cm,ybar,bar width=14pt,ylabel={coefficient d'information de test},
  symbolic x coords={logistic,CNN,LSTM,TCN,attention},xtick=data,
  xticklabels={{logistique sur\\variables},CNN,LSTM,TCN,attention},xticklabel style={align=center},tick label style={font=\small},
  label style={font=\small},ymin=-0.2,ymax=0.7,grid=major,grid style={gray!25},enlarge x limits=0.14,
  table/col sep=comma]
\addplot[fill=omDef!55,draw=omDef,error bars/.cd,y dir=both,y explicit]
  table[x=model,y=ic,y error=sd]{figdata/ml/08-sequence-models-on-order-books/models.csv};
\end{axis}
\end{tikzpicture}

\omcaption{Coefficient d'information de $\P(\text{hausse}) - \P(\text{baisse})$ avec la variation future, moyenné sur huit
sessions de test (et deux graines pour les modèles séquentiels)~; les barres montrent l'écart type entre sessions et graines. Le
\omterm{def:ml:linear-and-regularised-baselines:logistic}{régression logistique} lit dix variables construites à la main à l'instant de décision~; les modèles séquentiels lisent le carnet brut à dix niveaux
sur dix secondes. Données~: \texttt{ml\_lobseq.table}.}\label{fig:ml:lobseq:models}
\end{omfigure}

\section{Ce que la littérature reproduit}

Sur les données du chapitre (la \cref{fig:ml:lobseq:models}), la \omterm{def:ml:linear-and-regularised-baselines:logistic}{régression logistique} multinomiale sur les variables construites à la main du
livre~8 (chapitre~14~; déséquilibre du flux d'ordres (OFI) sur trois fenêtres, déséquilibre des files, flux et mouvements passés) atteint une précision de
0.781 et un IC de 0.543. Parmi les modèles séquentiels, le \omterm{def:ml:sequence-models-on-order-books:attention}{transformer} fait le mieux (précision 0.775, IC 0.426), puis le TCN
(0.766, 0.366), le LSTM (0.729, 0.292) et le CNN (0.727, 0.060)~: le CNN moyenne ses filtres sur la fenêtre et donc ne
peut distinguer le dernier instantané du premier, et il fait à peine mieux que de toujours prédire «~stable~». La dispersion entre sessions
et graines (0.15 à 0.20 d'IC pour les modèles séquentiels, 0.09 pour la \omterm{def:ml:linear-and-regularised-baselines:logistic}{régression logistique}) est aussi grande que les écarts.

Les modèles ont été soumis à un \omterm{def:ml:trees-and-boosting:early}{arrêt précoce} (early stopping) après une à quatre \omterm{def:ml:neural-networks-for-noisy-tabular-data:backprop}{époques}~: huit sessions représentent 8\,808 fenêtres, et un modèle
à plusieurs milliers de paramètres qui lit 800 nombres par fenêtre s'ajuste rapidement à leur bruit. Doubler les données aide (la \cref{fig:ml:lobseq:scaling})~:
l'IC du TCN passe de 0.13 sur quatre sessions à 0.30 sur huit et 0.44 sur seize, tandis que la \omterm{def:ml:linear-and-regularised-baselines:logistic}{régression logistique} reste
entre 0.54 et 0.57 tout du long. Les modèles séquentiels apprennent à partir du carnet brut ce que les variables construites à la main donnent
directement au modèle linéaire, et ils sont encore en train de l'apprendre avec seize sessions.

\begin{omfigure}
\begin{tikzpicture}
\begin{axis}[width=9.5cm,height=5.2cm,xlabel={sessions d'entraînement (dix minutes chacune)},ylabel={IC de test},
  tick label style={font=\small},label style={font=\small},xmin=2,xmax=18,ymin=0,ymax=0.7,xtick={4,8,16},
  grid=major,grid style={gray!25},legend columns=2,legend style={font=\footnotesize,at={(0.5,-0.3)},anchor=north,
  draw=none,fill=none,/tikz/every even column/.append style={column sep=8pt}},table/col sep=comma]
\addplot[omDef,very thick,mark=*] table[x=sessions,y=tcn]{figdata/ml/08-sequence-models-on-order-books/scaling.csv};
\addplot[omThm,very thick,mark=square*] table[x=sessions,y=logistic]{figdata/ml/08-sequence-models-on-order-books/scaling.csv};
\legend{TCN sur le carnet brut,régression logistique sur variables construites à la main}
\end{axis}
\end{tikzpicture}

\omcaption{IC de test à mesure que l'\omterm{def:ml:why-financial-machine-learning-is-different:sets}{ensemble d'entraînement} croît de quatre à seize sessions (une graine~; les mêmes huit sessions de test). Données~:
\texttt{ml\_lobseq.scaling}.}\label{fig:ml:lobseq:scaling}
\end{omfigure}

La littérature publiée va dans le même sens dès que l'évaluation est rigoureuse. DeepLOB, de Zhang, Zohren et Roberts, a surpassé
les algorithmes antérieurs sur l'indice de référence FI-2010 (Ntakaris et al.) et a conservé une précision hors échantillon (out-of-sample) stable sur une année de
cotations de la London Stock Exchange, y compris sur des instruments sur lesquels il n'avait pas été entraîné. Sirignano et Cont, en s'entraînant sur
des milliards de cotations américaines, ont mis en évidence une relation stable et universelle entre l'historique du flux d'ordres et la direction du
mouvement suivant, et qu'un modèle agrégé sur l'ensemble des actions bat les modèles propres à chaque action~: c'est l'échelle qui a rendu l'entrée brute exploitable. Kolm,
Turiel et Westray, sur 115 actions du Nasdaq, ont montré que des réseaux plus simples entraînés sur le flux d'ordres surpassent nettement
la plupart des modèles entraînés directement sur les carnets d'ordres. L'indice de référence LOBCAST, qui regroupe quinze modèles profonds publiés (Prata et al., 2024),
constate que tous perdent une part importante de leur performance sur des données sur lesquelles ils n'ont pas été calibrés. Deux leçons valent pour tout
modèle sur carnet d'ordres~: construire les variables stationnaires que la littérature sur la microstructure connaît déjà (flux d'ordres, déséquilibre,
livre~7, chapitre~8), et juger un modèle sur des sessions, des journées et des instruments qu'il n'a jamais vus.

\begin{method}[Entraîner un modèle séquentiel sur des carnets d'ordres]\label{met:ml:lobseq:method}
\begin{enumerate}
\item Découper par session (journée, instrument), jamais par fenêtre~; ajuster toute normalisation sur les sessions d'entraînement.
\item Fixer l'horizon de l'étiquette et la bande morte à partir de la transaction que le modèle sert~; indiquer la part de chaque classe et un score qui
ne récompense pas la prédiction de la classe majoritaire (l'IC, ou un P\&L tenant compte des coûts).
\item Partir d'un modèle linéaire sur les variables de flux d'ordres et de déséquilibre~; donner au modèle séquentiel les mêmes variables comme
canaux supplémentaires avant de lui demander de les retrouver.
\item Indiquer la dispersion entre sessions et graines, ainsi qu'une courbe d'apprentissage en sessions.
\end{enumerate}
\end{method}

\section{Tutoriel~: DeepLOB sur un ordinateur portable}\label{tut:ml:sequence-models-on-order-books:tut}

\begin{tutorial}
\textbf{Objectif.} Convertir des sessions simulées en fenêtres de carnet d'ordres, entraîner quatre modèles séquentiels et une \omterm{def:ml:linear-and-regularised-baselines:logistic}{régression logistique} de référence, et
mesurer la courbe d'apprentissage en sessions. \textbf{État final~:} les \cref{fig:ml:lobseq:models,fig:ml:lobseq:scaling}.
\begin{tutsteps}
\item \textbf{Fenêtres et étiquettes}, à l'intérieur de chaque session.
\omcode{code/firm/lobseq/firm_lobseq.py}{57}{73}{Fenêtres de snapshots, étiquettes de direction lissées et normaliseur ajusté sur les fenêtres d'entraînement.}\label{lst:ml:lobseq:windows}
\item \textbf{Le TCN et le modèle d'attention.}
\omcode{code/firm/lobseq/firm_lobseq.py}{97}{136}{Convolutions causales dilatées et transformer à une couche.}\label{lst:ml:lobseq:models}
\item \textbf{Exécuter} \texttt{ml\_lobseq.table()}, \texttt{scaling(4)}, \texttt{scaling(8)}, \texttt{scaling(16)} et
\texttt{fig\_lobseq.py}.
\end{tutsteps}
\textbf{Pour aller plus loin.} Donner au TCN les variables construites à la main comme dix canaux supplémentaires et voir quelle part de l'écart avec la
\omterm{def:ml:linear-and-regularised-baselines:logistic}{régression logistique} se referme~; remplacer l'étiquette lissée par le signe du prochain changement du mid.
\end{tutorial}

\section{Construction~: jeux de données et modèles séquentiels sur carnets d'ordres}\label{bld:ml:sequence-models-on-order-books:lobseq}

\begin{build}
\bfield{Rôle} Modèles séquentiels sur carnets d'ordres entraînés et évalués de la même façon sur des sessions simulées aujourd'hui, et sur les
bandes enregistrées du simulateur de bourse du livre~10 lorsqu'elles seront disponibles.
\bfield{Interface} \texttt{snapshots(tape, step, levels, start)}, \texttt{windows(S, mid, T, horizon, threshold)},
\texttt{Scaler}, \texttt{CNN}, \texttt{LSTMNet}, \texttt{TCN}, \texttt{Attention},
\texttt{train(model, X, y, Xv, yv, epochs, lr, batch, patience, seed)}, \texttt{evaluate(model, X, y, fwd)}.
\bfield{Règles} Les fenêtres d'une session restent dans un même découpage~; le carnet est reconstruit à partir des messages, sans aucun message
postérieur à l'instant du snapshot~; un seul fil CPU, algorithmes déterministes.
\bfield{Tests d'acceptation} \texttt{code/firm/lobseq/tests/}~: les snapshots correspondent à un carnet reconstruit à la main à un instant donné et
n'utilisent aucun message postérieur~; les fenêtres et les étiquettes sont vérifiées à la main sur une courte série~; la sortie du TCN à un pas ne dépend pas des
entrées ultérieures (causalité)~; chaque modèle s'entraîne sur un motif implanté et fait mieux que le hasard.
\bfield{Pour aller plus loin} Entrées de type market-by-order (âge et position dans la file de chaque ordre)~; un modèle agrégé sur plusieurs
instruments simulés, dans l'esprit de Sirignano et Cont.
\end{build}

\begin{omsources}
\item Z.~Zhang, S.~Zohren et S.~Roberts, «~DeepLOB: deep convolutional neural networks for limit order books~»,
\emph{IEEE Transactions on Signal Processing} 67(11), 2019.
\item J.~Sirignano et R.~Cont, «~Universal features of price formation in financial markets: perspectives from deep
learning~», \emph{Quantitative Finance} 19(9), 2019.
\item P.~N.~Kolm, J.~Turiel et N.~Westray, «~Deep order flow imbalance: extracting alpha at multiple horizons from the
limit order book~», \emph{Mathematical Finance} 33(4), 2023.
\item M.~Prata et al., «~LOB-based deep learning models for stock price trend prediction: a benchmark study~»,
\emph{Artificial Intelligence Review} 57, 2024.
\item A.~Ntakaris, M.~Magris, J.~Kanniainen, M.~Gabbouj et A.~Iosifidis, «~Benchmark dataset for mid-price forecasting of
limit order book data with machine learning methods~», \emph{Journal of Forecasting} 37(8), 2018.
\item S.~Hochreiter et J.~Schmidhuber, «~Long short-term memory~», \emph{Neural Computation} 9(8), 1997~; S.~Bai, J.~Z.~Kolter
et V.~Koltun, «~An empirical evaluation of generic convolutional and recurrent networks for sequence modeling~», arXiv,
2018~; A.~Vaswani et al., «~Attention is all you need~», \emph{NeurIPS}, 2017.
\end{omsources}

\section{Exercices}

\begin{exercise}[$\star$]\label{exo:ml:sequence-models-on-order-books:1}
Calculer le champ réceptif, en snapshots, d'un TCN dont les filtres ont une largeur de trois et des dilatations 1, 2, 4 et 8, puis avec une largeur de cinq.
\end{exercise}

\begin{exercise}[$\star$]\label{exo:ml:sequence-models-on-order-books:2}
Les fenêtres de test sont à 72.7~\% stables, à 14.7~\% à la baisse et à 12.5~\% à la hausse. Calculer la précision qu'atteint la prédiction systématique de «~stable~», et
celle d'un tirage au hasard selon les proportions des classes~?
\end{exercise}

\begin{exercise}[$\star$]\label{exo:ml:sequence-models-on-order-books:3}
Pour une tête d'attention sur $T = 20$ positions de largeur $h = 32$, calculer le nombre de multiplications qu'exige $\mathsf Q\mathsf K^\top$
et indiquer comment ce nombre croît avec $T$.
\end{exercise}

\begin{exercise}[$\star\star$]\label{exo:ml:sequence-models-on-order-books:4}
Pourquoi le CNN, qui moyenne sur la fenêtre, réussit-il si mal, et quelle petite modification y remédierait~?
\end{exercise}

\begin{exercise}[$\star\star$]\label{exo:ml:sequence-models-on-order-books:5}
L'étiquette moyenne le mid sur les dix secondes suivantes. Pourquoi cela la rend-elle plus facile à prédire que la variation suivante du mid,
et que signifie-t-elle pour le trading~?
\end{exercise}

\begin{exercise}[$\star\star$]\label{exo:ml:sequence-models-on-order-books:6}
\emph{Trouver la faille.} «~Nous avons découpé nos 50 jours de fenêtres de carnet d'ordres en \omterm{def:ml:why-financial-machine-learning-is-different:sets}{ensembles d'entraînement} et de test au hasard, et notre réseau atteint 88~\%
de précision.~»
\end{exercise}

\begin{exercise}[$\star\star\star$]\label{exo:ml:sequence-models-on-order-books:7}
\emph{Programmation.} Entraîner le LSTM sur seize sessions (\texttt{scaling(16, name='LSTM')}). Rapporter son IC et le comparer à celui du
TCN sur seize sessions.
\end{exercise}

\begin{exercise}[$\star\star\star$]\label{exo:ml:sequence-models-on-order-books:8}
Montrer qu'une pile de \omterm{def:ml:sequence-models-on-order-books:tcn}{convolutions causales} de largeur $k$ avec des dilatations $1, 2,\dots,2^{L-1}$ possède un champ réceptif de
$1 + (k - 1)(2^L - 1)$ positions et que sa sortie en $t$ ne dépend d'aucune entrée postérieure à $t$.
\end{exercise}

\section{Problème~: DeepLOB sur un ordinateur portable}

\begin{problem}[{Problème du week-end --- le carnet brut contre les variables construites à la main}]\label{pb:ml:sequence-models-on-order-books:1}
Les sessions du chapitre~: huit pour l'entraînement, trois pour l'arrêt, huit pour le test, huit autres pour agrandir l'\omterm{def:ml:why-financial-machine-learning-is-different:sets}{ensemble d'entraînement}.

\medskip\noindent\textbf{Partie I --- Les données.}
\begin{enumerate}
\item Que contient un instantané et une fenêtre~?
\item Comment l'étiquette est-elle construite, et quelles sont les proportions des classes~?
\item Pourquoi la précision est-elle un mauvais score ici, et que retient le chapitre à la place~?
\item Pourquoi les fenêtres ne doivent-elles jamais chevaucher la séparation entre sessions~?
\end{enumerate}

\medskip\noindent\textbf{Partie II --- Les modèles.}
\begin{enumerate}[resume]
\item Quels sont les champs réceptifs des quatre modèles de séquences sur la fenêtre~?
\item Que vaut chaque modèle (précision et IC)~?
\item Que vaut la \omterm{def:ml:linear-and-regularised-baselines:logistic}{régression logistique} sur les variables construites à la main~?
\item Quelle est l'amplitude de l'écart entre les sessions et les graines~?
\end{enumerate}

\medskip\noindent\textbf{Partie III --- Les données (suite).}
\begin{enumerate}[resume]
\item Au bout de combien d'\omterm{def:ml:neural-networks-for-noisy-tabular-data:backprop}{époques} les modèles s'arrêtent-ils, et pourquoi~?
\item Que vaut le TCN avec quatre, huit et seize sessions d'entraînement~?
\item La \omterm{def:ml:linear-and-regularised-baselines:logistic}{régression logistique} profite-t-elle de davantage de sessions~? Pourquoi pas~?
\item Que faudrait-il que la courbe fasse pour que le TCN dépasse la référence~?
\end{enumerate}

\medskip\noindent\textbf{Partie IV --- Le verdict.}
\begin{enumerate}[resume]
\item Qu'ont trouvé Kolm, Turiel et Westray sur les carnets bruts face au flux d'ordres~?
\item Qu'a trouvé l'indice de référence LOBCAST sur les modèles publiés face à de nouvelles données~?
\item Qu'est-ce qui a permis au modèle à entrées brutes de Sirignano et Cont de fonctionner~?
\item Énoncer le \emph{résultat nommé}~: l'IC et la précision du meilleur modèle de séquences face à la \omterm{def:ml:linear-and-regularised-baselines:logistic}{régression logistique} sur le déséquilibre du flux d'ordres,
avec l'écart entre les graines et les sessions.
\item Que mettre en production pour une prévision à dix secondes, et que tenter ensuite~?
\item Comment \omterm{def:ml:sequence-models-on-order-books:attention}{transformer} l'IC en décision de trading, étant donné les étiquettes du chapitre~2~?
\item Qu'est-ce qui change dans le simulateur du livre~10 avec plusieurs instruments~?
\item En une phrase~: que faut-il à un modèle de séquences pour battre une bonne variable~?
\end{enumerate}
\end{problem}

\section{Questions d'entretien}

\begin{interviewq}[$\star$ \iqroles{ingénieur ML, chercheur}]\label{iq:ml:sequence-models-on-order-books:1}
Qu'est-ce qu'une \omterm{def:ml:sequence-models-on-order-books:tcn}{convolution causale}, et pourquoi importe-t-elle pour un modèle de trading~?
\end{interviewq}

\begin{interviewq}[$\star\star$ \iqroles{ingénieur ML}]\label{iq:ml:sequence-models-on-order-books:2}
Expliquez l'auto-attention en quelques équations. Que coûte-t-elle en fonction de la longueur de la séquence~?
\end{interviewq}

\begin{interviewq}[$\star\star$ \iqroles{chercheur}]\label{iq:ml:sequence-models-on-order-books:3}
Votre modèle de carnet d'ordres a une précision de 80~\% sur une étiquette à trois classes fondée sur le prix milieu. Quelles questions posez-vous avant d'être
impressionné~?
\end{interviewq}

\begin{interviewq}[$\star\star$ \iqroles{chercheur, trader}]\label{iq:ml:sequence-models-on-order-books:4}
Des niveaux bruts du carnet ou des variables de flux d'ordres comme entrées~? Argumentez à partir des faits.
\end{interviewq}

\begin{interviewq}[$\star\star$ \iqroles{ingénieur ML}]\label{iq:ml:sequence-models-on-order-books:5}
Comment découpez-vous les données de carnet d'ordres en \omterm{def:ml:why-financial-machine-learning-is-different:sets}{ensembles d'entraînement} et de test, et pourquoi~?
\end{interviewq}

\begin{interviewq}[$\star\star\star$ \iqroles{chercheur, ingénieur ML}]\label{iq:ml:sequence-models-on-order-books:6}
Pourquoi un LSTM pourrait-il sous-performer un TCN sur de courtes fenêtres de données de carnet d'ordres, et quand attendriez-vous l'inverse~?
\end{interviewq}
